\chapter{序論}
本章では，本研究の研究背景と目的について述べる．

% \section{研究背景}
% 近年，都市や地域社会における情報基盤の構築が重要視されている．
% これまで，センシングデータや地方自治体が公開する統計データを活用し，まちの情報基盤に関する研究を行ってきた\cite{cds41_mine_proposal}．
% これらのデータは客観的でまちの状況を把握する上で重要である一方，いくつかの課題が存在する．
% まず，統計データは更新頻度が低く，リアルタイム性に欠けるという問題がある．
% また，センサデータは設置地点での定点的な情報に限られ，広範囲な状況把握には多数のセンサが必要となるため，設置コストや維持管理の負担が大きい．
% さらに重要な点として，これらの定量的データだけでは市民の主観的な意見や感情的な反応を捉えることができないという限界がある．
% このような主観的・情動的要素を考慮しない都市分析は、政策決定や都市サービスの設計において、市民の実感と乖離した判断を導く可能性がある。

% まちの状況を把握するための手段として従来の研究では，X（旧Twitter）などのSNSコメントが注目されてきた\cite{cds41_twitter_disaster,cds41_twitter_event}．
% 主に災害時における市民の反応を捉える情報源として活用された例が多い．このような背景から，市民の主観的・情動的な反応が即時的かつ大規模に集積しているSNSコメントを新たな「ヒューマンセンサ」\cite{human_sensing}として捉え活用する．
% 地震や台風、津波、内水氾濫などを代表とする災害時には，市民の反応がSNS上に多数投稿されるため，従来のセンサデータでは捉えきれない現場の状況や市民の困りごとや行動を把握できる可能性がある．
% これにより，まちの状況把握や災害時対策の高度化が期待される．

% その上で，SNSコメントを情報源として活用する際には，固有の課題も存在する．
% 研究で利用されるSNSの代表であるXに注目すると，近年の仕様変更により無料枠でのデータ収集量に制限が設けられるなど，大規模なデータ収集が困難になりつつある．
% また，SNSコメントにはノイズや誤情報が含まれるという本質的な課題もあり，有用な情報を抽出するためには適切な処理が必要となる．
% SNSから収集したデータは、都市の状況把握を目的とした分析に適用するには未加工の状態では不十分であり、分析など別システムへの応用も含めると，構造化やノイズ除去などの前処理を施すことが求められ、継続的な蓄積，補完，再構成が必要となる．

% これらの課題から、まちの人々の状況を把握するための情報源としてSNSの活用が有力であり、従来使用されてきたX以外のSNSプラットフォームの検討も重要となっている．
% そして、SNSデータのソースを検討することや、そこからまちの状況を把握する定性、定量的に分析を可能とする基盤が必要である．

% \section{研究背景}
% 近年，都市や地域社会における持続的かつ高度な運営を実現するため，都市の状況を的確に把握する情報基盤の構築が重要視されている．
% これまでに，センシングデータや地方自治体が公開する統計データを活用し，都市の状況把握を目的とした情報基盤に関する研究が行われてきた\cite{cds41_mine_proposal}．
% これらのデータは客観性が高く，都市の状態を定量的に把握する上で有用である一方，いくつかの課題が指摘されている．

% まず，統計データは更新頻度が低く，リアルタイムな状況変化を捉えることが困難である．
% また，センサデータは設置地点に依存した定点観測に限られるため，広域的な状況把握には多数のセンサが必要となり，設置および維持管理に伴うコスト負担が大きい．
% さらに，これらの定量的データのみでは，市民の主観的な意見や情動的な反応といった，人間の内面的な状態を十分に捉えることができないという限界がある．
% このような主観的・情動的要素を考慮しない都市分析は，政策決定や都市サービス設計において，市民の実感と乖離した判断を導く可能性がある．

% こうした課題に対し，都市の状況を把握する新たな情報源として，X（旧Twitter）をはじめとするSNS上に投稿されるコメントが注目されてきた\cite{cds41_twitter_disaster,cds41_twitter_event}．
% 特に災害発生時には，市民の行動や感情，現場の状況に関する情報が即時的かつ大量に投稿されることから，従来のセンサデータでは把握が困難であった状況を補完する情報源として活用されている．
% このような背景から，SNSコメントを市民の主観的・情動的反応が集積した「ヒューマンセンサ」\cite{human_sensing}として捉え，都市分析に活用する試みが進められている．
% 地震，台風，津波，内水氾濫などの災害時においては，SNSを通じて市民の困りごとや行動，現場の混乱状況が共有されるため，都市の状況把握や災害対応の高度化に資する可能性がある．

% 一方で，SNSコメントを情報源として活用する際には，固有の課題も存在する． 代表的なSNSであるXにおいては，近年の仕様変更により，無料枠で取得可能なデータ量が制限されるなど，大規模かつ継続的なデータ収集が困難になりつつある．
% また，SNSコメントにはノイズや誤情報が含まれるという特性があり，有用な情報を抽出するためには適切な前処理が不可欠である．
% 都市の状況把握を目的とした分析にSNSデータを適用するためには，収集したデータをそのまま利用するのではなく，構造化やノイズ除去などの前処理を施し，継続的な蓄積，補完，再構成を行う必要がある．

% 以上の背景から，都市における人々の状況を把握する情報源としてSNSの活用は有力であり，従来主に利用されてきたX以外のSNSプラットフォームについても検討することが重要である．
% さらに，SNSデータを対象とし，都市の状況を長期的にさまざまな角度で分析を可能とする情報基盤の構築が必要である．

\section{研究背景}
近年，都市や地域社会における持続的かつ高度な運営を実現するため，都市の状況を的確に把握する情報基盤の構築が重要視されている．
これまで，センシングデータや地方自治体が公開する統計データを活用し，都市の状況把握を目的とした情報基盤に関する研究が行われてきた\cite{cds41_mine_proposal}．
これらのデータは客観性が高く，都市の状態を定量的に把握する上で重要である一方，いくつかの課題が指摘されている．

まず，統計データは更新頻度が低く，リアルタイムな状況変化を捉えることが困難である．
また，センサデータは設置地点に依存した定点観測に限られるため，広域的な状況把握には多数のセンサが必要となり，設置および維持管理に伴うコスト負担が大きい．
さらに，これらの定量的データのみでは，市民の主観的な意見や情動的な反応といった，人間の内面的な状態を十分に捉えることができないという限界がある．
このような主観的・情動的要素を考慮しない都市分析は，政策決定や都市サービス設計において，市民の実感と乖離した判断を導く可能性がある．

こうした課題に対し，都市の状況を把握する新たな情報源として，X（旧Twitter）をはじめとするSNS上に投稿されるコメントが注目されてきた\cite{cds41_twitter_disaster,cds41_twitter_event}．
特に災害発生時には，被災者の感情や住民ニーズ，現場の状況に関する情報が即時的かつ大量に投稿されることから，従来のセンサデータでは把握が困難であった状況を補完する情報源として有効であることが示されている．
このような背景から，市民の主観的・情動的な反応が集積するSNSコメントを新たな「ヒューマンセンサ」\cite{human_sensing}として捉え，都市分析や災害対応に活用する研究が進められてきた．

一方で，既存研究の多くは災害発生後の限定的な期間を対象とした分析にとどまっており，平常時を含めた長期的・継続的な分析基盤としての活用には至っていない．
また，近年，ソーシャルメディアデータを活用した社会動向分析の研究基盤そのものが，大きな転換点を迎えている．
代表的なSNSプラットフォームであるXにおいては，従来は比較的自由に利用可能であったAPIへのアクセスが厳しく制限され，有償化が進められている．
これにより，大規模かつ網羅的なデータを継続的に収集することが著しく困難となり，従来のSNSデータに依存した研究アプローチの持続可能性が低下している．

さらに，SNS上には偽情報や誤情報，自動化されたアカウントによる投稿といったノイズが大量に含まれている．
これらのデータから信頼性の高い知見を抽出するためには，高度で複雑なフィルタリング処理が不可欠であり，ノイズの傾向や特性を把握するためには，十分な量のデータ蓄積が求められる．
このような量的・質的制約は，SNS上のテキストデータを主たる情報源とする都市分析や災害分析の信頼性に大きな影響を及ぼしている．
そのため，SNSデータは未加工の状態では都市の状況把握に適用することが難しく，構造化やノイズ除去といった前処理を施した上で，継続的な蓄積，補完，再構成を行う分析基盤が必要となる．

以上の背景から，都市における人々の状況を把握する情報源としてSNSの活用は有力であるものの，主に利用されてきたX以外のSNSプラットフォームについても検討することが重要である．
その上で，SNSデータを対象とし，平常時から災害時に至るまで都市の状況を多角的に分析可能とする情報基盤が必要である．

% \section{研究目的}
% 本研究の目的は，災害時における市民の反応を捉える新たな情報源として，YouTubeライブ配信コメントを活用する分析基盤を構築し，その特徴と有用性を明らかにすることである．

% 具体的には，以下の3点を達成することを目指す．
% 第一に，YouTubeライブ配信コメントの収集・蓄積・分析を行うシステムを構築する．
% これにより，災害時のコメントを継続的に収集し，体系的な分析を可能にする．
% 第二に，既存のSNSであるXとYouTubeライブ配信コメントの仕様を比較し，データ収集量やコメント特性の観点から両者の違いを明確にする．
% 第三に，収集したコメントに対して，時系列分析，クラスタリング，感情分析，ユーザ分類などの多角的な分析を実施し，YouTubeライブチャットが災害時の情報源として持つ特徴を明らかにする．

% これらを通じて，まちの定量的なデータと市民の主観的な反応を接続し，災害対応支援の高度化を可能とする新たな情報基盤の構築に貢献する．

\section{研究目的}
本研究では，市民の反応が現れやすい災害時に注目する。
まちの状況把握や災害時対策への応用を目的として，SNSコメントを活用した持続的な分析基盤の構築と，新たなデータ収集ソースの有効性を明らかにすることを目的とする．
特に，従来の研究で多用されてきたX（旧Twitter）に加え，新たなSNSコメントの収集ソースとしてYouTubeライブ配信のコメントに着目する．

本研究では，以下の3つの研究課題を設定する．
\begin{enumerate}
    \item YouTubeライブ配信コメントの特徴を明らかにし，Xに投稿されるポストとの比較を行う．

    \item YouTubeに投稿されるコメントを長期的に収集し，平常時から災害時まで継続的な分析を可能とするデータ収集・分析基盤を提案する．

    \item 地震発生時を対象として，コメント数の時系列変化やピーク検出，感情分析モデルの適用を行い，収集したコメントの特性を定量的に明らかにする．
\end{enumerate}

これらを通じて，YouTubeライブ配信コメントが都市の状況把握および災害対応支援において有効な情報源となり得るかを検証し，SNSを活用した都市分析基盤の高度化に資する知見を提供する．